\section{Static correlation: H$_2$ dissociation and the weak-interaction limit (Figure 2)}\label{sec:static}

\begin{reviewclaim}
``\dots dynamic correlation is weak compared to static correlation effects, i.e.\ correlation between
configurations in the multireference wavefunction \dots\ The weaker the interaction between the two
nonbonding molecular orbitals, the stronger the static correlation (Figure 2). This is clearly
illustrated by \dots\ the H$_2$ homolytic dissociation curve or the singlet ground state of square
cyclobutadiene \dots'' (p.~442)
\end{reviewclaim}

\subsection{H$_2$ in a minimal basis: exact two-configuration treatment}
Basis: one $1s$ function on each nucleus, $\chi_A,\chi_B$, overlap $S$. By symmetry the MOs are
\begin{equation}
  \sigma_g\equiv\varphi_1=\frac{\chi_A+\chi_B}{\sqrt{2(1+S)}},\qquad
  \sigma_u\equiv\varphi_2=\frac{\chi_A-\chi_B}{\sqrt{2(1-S)}} .
\end{equation}
$\varphi_1$ is even and $\varphi_2$ odd under inversion, so $h_{12}=\eri1112=\eri2221=0$ and \cref{eq:Hsinglet}
splits. The OSS ($^1\Sigma_u^+$) is alone, and the two closed shells ($^1\Sigma_g^+$ symmetry) mix:
\begin{equation}\label{eq:h2ci}
  \mathbf H(^1\Sigma_g^+)=\begin{pmatrix}2h_{11}+J_{11}&K_{12}\\K_{12}&2h_{22}+J_{22}\end{pmatrix},\qquad
  E_{\text{CI}}=\tfrac12\Bigl[E_1+E_2-\sqrt{(E_2-E_1)^2+4K_{12}^2}\Bigr],
\end{equation}
with $E_1=2h_{11}+J_{11}$ (the RHF energy), $E_2=2h_{22}+J_{22}$, and
$E(^3\Sigma_u^+)=h_{11}+h_{22}+J_{12}-K_{12}$, $E(^1\Sigma_u^+)=h_{11}+h_{22}+J_{12}+K_{12}$
(nuclear repulsion $1/R$ added to all).

\subsection{The dissociation limit, step by step}
Let $R\to\infty$, so $S\to0$ and all two-centre overlap densities $\chi_A\chi_B\to0$. Define the on-site
repulsion $U=\eri AAAA$, the inter-site repulsion $V=\eri AABB\to1/R$, and the H-atom energy $E_H$.
\begin{enumerate}
  \item Orbital products: $\varphi_1^2=\tfrac12(\chi_A^2+\chi_B^2+2\chi_A\chi_B)\to\tfrac12(\chi_A^2+\chi_B^2)$,
        $\varphi_2^2\to\tfrac12(\chi_A^2+\chi_B^2)$, $\varphi_1\varphi_2\to\tfrac12(\chi_A^2-\chi_B^2)$.
  \item Two-electron integrals (each product contributes $\tfrac12\cdot\tfrac12$):
        \begin{align*}
        J_{11}=J_{22}=J_{12}&\to\tfrac14\bigl[\eri AAAA+\eri AABB+\eri BBAA+\eri BBBB\bigr]=\tfrac12(U+V),\\
        K_{12}=\eri1212&\to\tfrac14\bigl[\eri AAAA-\eri AABB-\eri BBAA+\eri BBBB\bigr]=\tfrac12(U-V).
        \end{align*}
  \item One-electron integrals: $h_{11}=h_{22}\to h_{AA}=E_H-\tfrac1R$ (an electron on A also feels nucleus B
        as a point charge at distance $R$).
  \item RHF: $E_{\text{RHF}}=2h_{11}+J_{11}+\tfrac1R=2E_H-\tfrac2R+\tfrac12(U+V)+\tfrac1R\to2E_H+\tfrac U2-\tfrac1{2R}$.
        \emph{The RHF error at dissociation is $U/2$}, which does not vanish.
  \item Triplet: $E_T=2h_{11}+J_{12}-K_{12}+\tfrac1R=2E_H-\tfrac2R+V+\tfrac1R\to2E_H$. Correct.
  \item CI: the diagonal elements are now equal ($E_1=E_2$), so by \cref{eq:h2ci}
        $E_{\text{CI}}=E_1-K_{12}=2E_H+\tfrac U2-\tfrac1{2R}-\tfrac12(U-V)\to2E_H$. Correct.
        The eigenvector has $c_1=-c_2=\tfrac1{\sqrt2}$: \emph{two configurations with equal weight}. This is
        maximal static correlation.
  \item What that eigenvector is: $\varphi_1(1)\varphi_1(2)-\varphi_2(1)\varphi_2(2)\to\tfrac12[(\chi_A+\chi_B)(\chi_A+\chi_B)-(\chi_A-\chi_B)(\chi_A-\chi_B)]
        =\chi_A(1)\chi_B(2)+\chi_B(1)\chi_A(2)$, the purely covalent Heitler--London function. By contrast
        $\varphi_1(1)\varphi_1(2)$ alone is $\tfrac12[\text{covalent}+\text{ionic}]$, which is the origin of the
        $U/2$ error in step 4.
  \item OSS: $E(^1\Sigma_u^+)=E_T+2K_{12}\to2E_H+U-V$. It is the antisymmetric \emph{ionic} combination.
\end{enumerate}

\begin{verifybox}
T5 (\texttt{t03\_t05\_hund\_virial\_h2.py}), STO-3G, PySCF vs.\ the closed forms at $R=1.4,3,6,20$ bohr.
RHF, 2$\times$2 CI\,=\,FCI and ROHF triplet agree with \cref{eq:h2ci} to $4\times10^{-16}\,E_h$ (max.\ deviation, tight SCF). At $R=20$:
$E_{\text{CI}}=E_T=-0.933164=2E_H$; $E_{\text{RHF}}-2E_H=0.3623=K_{12}$; $J_{11}-K_{12}=0.050=1/R$ (so $U=J_{11}+K_{12}=0.7746$);
$E(^1\Sigma_u^+)-2E_H=0.7246=U-1/R$. Every limit in steps 4--8 is reproduced. \textbf{PASS}.
\end{verifybox}

\subsection{General weak-interaction limit (Figure 2b, twisted ethylene)}
Take the $2e/2o$ model with orbitals of different symmetry ($h_{ab}=\eri aaab=\eri bbba=0$), orbital
energy gap $\Delta$ ($h_a=-\Delta/2$, $h_b=+\Delta/2$), and for simplicity $J_{aa}=J_{bb}\equiv J_d$,
$J_{ab}\equiv J_x$. This describes the $\pi/\pi^*$ pair of ethylene, where $\Delta$ is the $\pi$--$\pi^*$
splitting, which goes to $0$ at a $90^\circ$ twist.
\begin{enumerate}
  \item Closed-shell block: $\begin{pmatrix}J_d-\Delta&K\\K&J_d+\Delta\end{pmatrix}$, eigenvalues $J_d\pm\sqrt{\Delta^2+K^2}$.
  \item Ground singlet: $E_S=J_d-\sqrt{\Delta^2+K^2}$, with coefficient ratio
        $\dfrac{c_2}{c_1}=\dfrac{E_S-(J_d-\Delta)}{K}=\dfrac{\Delta-\sqrt{\Delta^2+K^2}}{K}$.
  \item Limits: $\Delta\gg K$ gives $c_2/c_1\to0$ (one configuration, closed shell, no static correlation);
        $\Delta\to0$ gives $c_2/c_1\to-1$ (two configurations of equal weight, a pure diradical).
  \item Gap to the triplet ($E_T=J_x-K$): $E_S-E_T=(J_d-J_x)+K-\sqrt{\Delta^2+K^2}$.
\end{enumerate}
\begin{resultbox}
Static correlation (the weight of the second configuration) is controlled by $\Delta/K$. The smaller
the splitting of the two frontier orbitals (``weaker interaction''), the more multiconfigurational the
singlet becomes, and the closer it comes to the triplet. This is the content of Figure~2(a,b).
\end{resultbox}
\begin{verifybox}
T6 (\texttt{t06\_t09\_models.py}): eigenvalue, ratio and both limits verified symbolically. \textbf{PASS}.
Balkov\'a \& Bartlett, \emph{J.\ Chem.\ Phys.}\ \textbf{101}, 8972 (1994), doi:10.1063/1.468057
[abstract]: MR-CC places the square-cyclobutadiene singlet 6.9~kcal/mol below the triplet.
\end{verifybox}
Figure~2(c) (spin crossover in a transition-metal complex) is qualitative in the review and contains no
mathematics to derive.
